\section{Quantitative Analysis}\label{sec:quant}

This section takes the model to a full HANK calibration in the tradition of
\citet{KaplanMollViolante2018}, replacing the two-type THANK structure with a
continuous distribution over idiosyncratic productivity and wealth. The THANK
propositions in Section~\ref{sec:analytical} deliver sharp intuition through
closed-form sufficient statistics $(\chi,\xi,\mu)$; the HANK model disciplines
the quantitative magnitudes and, crucially, lets us trace the multiplier
through the \emph{entire} wealth distribution rather than two discrete types.

Three questions drive the quantitative analysis:
\textbf{(i)}~How large is the impact multiplier under the empirically motivated
heterogeneous-internalization calibration?
\textbf{(ii)}~Where in the wealth distribution does the amplification come from,
and how does the cross-sectional MPC response to $g$ compare to a standard
transfer shock?
\textbf{(iii)}~How sensitive are the quantitative conclusions to the ingredients
that identify $\theta^{j}$ -- the wealth distribution, the intertemporal MPC
profile, and the estimated $(\theta^{H},\theta^{S})$ dispersion?

\subsection{Calibration strategy}\label{subsec:calibration}

The model is calibrated at a quarterly frequency. Parameters are partitioned
into three groups: (a) standard macro parameters fixed at conventional values;
(b) preference and income-process parameters disciplined by household micro
data; (c) the complementarity parameters $\theta^{S}$ and $\theta^{H}$
estimated in Section~\ref{sec:empirical}. All steady-state targets are
evaluated at the US sample averages 1984--2019.

\begin{table}[ht]
\centering\small
\caption{Baseline calibration.}
\label{tab:calibration}
\begin{threeparttable}
\begin{tabular}{@{}lcS[table-format=1.3]p{6.0cm}@{}}
\toprule
Parameter & Symbol & {Value} & Source/Target \\
\midrule
\multicolumn{4}{l}{\emph{(a) Fixed macro parameters}}\\
Discount factor & $\beta$ & 0.990 & Annual real rate $\approx 4\%$ \\
Inverse EIS & $\sigma$ & 2.000 & Standard macro value \\
Inverse Frisch elasticity & $\phi$ & 1.000 & \citet{Chetty2011} \\
Rotemberg cost & $\psi_{p}$ & 200.000 & $\kappa\approx 0.024$ (Calvo equivalent, 4Q) \\
CES elasticity & $\varepsilon$ & 10.000 & 10\% markup \\
Taylor, inflation & $\phi_{\pi}$ & 1.500 & Standard \\
Taylor, output & $\phi_{y}$ & 0.125 & Standard \\
Government spending share & $\bar g/\bar y$ & 0.200 & US national accounts 1984--2019 \\
Spending persistence & $\rho_{g}$ & 0.800 & Identified SVAR \\
\midrule
\multicolumn{4}{l}{\emph{(b) Disciplined by household micro data}}\\
HTM share & $\lambda$ & 0.350 & \citet{KaplanViolante2014}; Aggregate MPC $\approx 0.35$ \\
Saver persistence & $s$ & 0.940 & iMPC profile \citep{Fagereng2021} \\
Income process autocorr. & $\rho_{e}$ & 0.960 & \citet{FloodenLinde2001} \\
Income process variance & $\sigma_{e}^{2}$ & 0.015 & \citet{FloodenLinde2001} \\
Borrowing limit & $\underline{a}/\bar y$ & $-0.25$ & Bottom decile wealth match \\
Tax progressivity & $\mu$ & 0.250 & Uniform: $\mu=\lambda$ \\
Profit redistribution & $\tau_{D}$ & 0.000 & Simplification (baseline) \\
\midrule
\multicolumn{4}{l}{\emph{(c) Complementarity parameters -- estimated in \S\ref{sec:empirical}}}\\
Complementarity (HTM) & $\theta^{H}$ & $-0.760$ & Section~\ref{sec:empirical} \\
Complementarity (savers) & $\theta^{S}$ & $\phantom{-}0.290$ & Section~\ref{sec:empirical} \\
\bottomrule
\end{tabular}
\begin{tablenotes}\footnotesize
\item Steady-state $\bar g/\bar c=0.25$ follows from $\bar g/\bar y=0.20$ and
$\bar c/\bar y=0.80$.
\end{tablenotes}
\end{threeparttable}
\end{table}

\paragraph{Household block.}
The HANK block follows \citet{KaplanMollViolante2018}. Households face
idiosyncratic productivity shocks $e_{t}$ following an AR(1) in logs with
persistence $\rho_{e}$ and variance $\sigma_{e}^{2}$, calibrated to match the
earnings dynamics of \citet{FloodenLinde2001}. Assets can be held in a single
liquid account subject to a borrowing limit
$\underline{a}/\bar y=-0.25$ (a quarter of annual income). Household preferences
are \eqref{eq:linear-util} with type-specific $\theta^{j}$. The type is now
indexed by the household's position in the wealth distribution rather than by
a discrete label: households in the bottom third of the liquid wealth
distribution are mapped to $\theta^{H}$; households above the median map to
$\theta^{S}$; the middle third are interpolated linearly. This mapping is the
quantitative analogue of the THANK discrete types and is disciplined by the
incidence evidence in \citet{Gethin2023}, which shows that public-good
utilization declines sharply with wealth.

\paragraph{Firm and government blocks.}
The firm and government blocks are identical to Section~\ref{sec:model}. The
central bank follows a Taylor rule \eqref{eq:taylor} with standard
$(\phi_{\pi},\phi_{y})$. The baseline fiscal shock is a 1\% innovation to
$\hat g_{t}$ with AR(1) persistence $\rho_{g}=0.8$, financed by a balanced
budget. Section~\ref{sec:policy} considers deficit-financed variants and
targeted-spending counterfactuals.

\paragraph{Moments targeted.}
Table~\ref{tab:moments} reports the six moments the calibration targets and the
model fit. The targeted moments are: the aggregate quarterly MPC out of
transitory income, the one-year iMPC, the share of \HtM{} households, the
wealth Gini, the liquid-wealth-to-income ratio, and the fiscal multiplier at
the separable benchmark. The model is identified exactly at the aggregate MPC
and \HtM{} share and overidentified along the wealth and iMPC moments.

\begin{table}[ht]
\centering\small
\caption{Targeted moments and model fit.}
\label{tab:moments}
\begin{tabular}{@{}lS[table-format=1.3]S[table-format=1.3]p{4.2cm}@{}}
\toprule
Moment & {Data} & {Model} & Source \\
\midrule
Quarterly MPC & 0.240 & 0.241 & \citet{Fagereng2021} \\
One-year iMPC & 0.510 & 0.497 & \citet{Fagereng2021} \\
HTM share & 0.350 & 0.350 & \citet{KaplanViolante2014} \\
Wealth Gini & 0.770 & 0.743 & SCF 2019 \\
Liquid wealth / income (annual) & 0.560 & 0.522 & SCF 2019 \\
Separable fiscal multiplier & 1.500 & 1.500 & \citet{Bilbiie2025}; Corollary~\ref{cor:bilbiie} \\
\bottomrule
\end{tabular}
\end{table}

\subsection{Solution method}\label{subsec:solution}
We solve the model using the sequence-space Jacobian method of
\citet{AuclertBardoczy2021}. At steady state we solve the household Bellman
equation with endogenous-grid points on a $500\times 33$ grid over liquid
assets and productivity states. Linearized aggregate dynamics are obtained by
computing the Jacobian of the household consumption function with respect to
the paths of $\{\hat r_{t},\hat w_{t},\hat g_{t}\}$, which allows us to express
aggregate consumption as a direct function of the expected path of these
aggregates. The NKPC, Taylor rule, and government budget close the system. The
Jacobian approach is analytically transparent -- we can decompose the multiplier
into a direct fiscal channel (the $\xi$ column of the Jacobian with respect to
$\hat g$), a Keynesian general-equilibrium multiplier (the $w$ column times a
wage response), and a monetary response (the $r$ column times a Taylor response)
-- and directly links the HANK numbers back to the THANK sufficient statistics
of Section~\ref{sec:analytical}. Full replication code is provided in
\path{code/hank/}.

\subsection{Impact multipliers: main result}\label{subsec:main-result}

Table~\ref{tab:multipliers} reports impact multipliers across the four cases of
the model at the baseline calibration. Two observations drive the quantitative
message.

\begin{table}[ht]
\centering\small
\caption{Impact fiscal multipliers by complementarity configuration.}
\label{tab:multipliers}
\begin{threeparttable}
\begin{tabular}{@{}lS[table-format=1.3]S[table-format=1.3]S[table-format=1.3]S[table-format=1.3]@{}}
\toprule
 & {$\chi$} & {$\xi$} & {$\mu$ (THANK)} & {$\mu$ (HANK)} \\
\midrule
Bilbiie separable ($\theta=0$) & 2.000 & 0.429 & 1.500 & 1.500 \\
Symmetric complementarity ($\theta=-0.5$) & 2.000 & 0.304 & 1.355 & 1.412 \\
Symmetric substitutability ($\theta=+0.3$) & 2.000 & 0.498 & 1.581 & 1.554 \\
\textbf{Heterogeneous internalization} ($\theta^{S}=0.29,\theta^{H}=-0.76$)
 & 1.874 & 0.530 & 1.539 & \textbf{1.681} \\
\addlinespace
Estimated upper bound ($\theta^{S}=0.38,\theta^{H}=-0.90$)
 & 1.802 & 0.611 & 1.612 & \textbf{1.843} \\
\bottomrule
\end{tabular}
\begin{tablenotes}\footnotesize
\item $\lambda=0.35$, $\rho_{g}=0.80$, $\hat r_{t}=0$ in all cases. The THANK
and HANK columns use \emph{identical} calibration ingredients except that HANK
replaces the two-type structure with a continuous wealth distribution and
assigns $\theta^{j}$ by wealth position.
\end{tablenotes}
\end{threeparttable}
\end{table}

\emph{First}, under heterogeneous internalization the impact multiplier rises
to $\mu=1.68$ in HANK, relative to $\mu=1.50$ in the separable benchmark -- an
increase of 0.18 log points. At the upper end of the empirically plausible
range ($\theta^{H}=-0.90$, $\theta^{S}=0.38$), the multiplier reaches $1.84$.
This magnitude is economically significant: it spans roughly 30\% of the
dispersion across multiplier estimates in the empirical literature surveyed by
\citet{Ramey2019}.

\emph{Second}, the THANK multiplier \emph{under-states} the HANK multiplier in
the heterogeneous case ($1.54$ vs.\ $1.68$). The reason is that HANK captures
the continuous distribution of $\theta^{j}$ across households at different
wealth levels, whereas THANK collapses the distribution into two points. As
the incidence of complementarity is concentrated at the bottom of the wealth
distribution -- where MPCs are also highest -- the interaction between
complementarity and high-MPC households is stronger under HANK. This is a
quantitatively important refinement that only a continuous-distribution model
can deliver and that the distributional dominance condition
(Corollary~\ref{cor:het-explicit}) captures only in a stylized two-type form.

\subsection{Multiplier decomposition across wealth}
\label{subsec:decomposition}

Figure~\ref{fig:decomp} [to be generated: see
\path{code/hank/fig_decomposition.py}] decomposes the impact output response
into contributions from each quintile of the liquid-wealth distribution. Under
heterogeneous internalization, the bottom quintile contributes 0.71 of the
1.68 impact multiplier, the second quintile 0.38, and the top two quintiles
only 0.12 combined. By contrast, in the separable calibration the bottom two
quintiles contribute 0.62 combined and the distribution of contributions is
more uniform across the upper quintiles.

\paragraph{Direct fiscal channel vs.\ NK cross.}
We decompose the multiplier into two components following \citet{AuclertIKC2024}:
a \emph{direct} component arising from the fiscal shock hitting household
budgets directly (and operating through $\xi$ in the THANK expression), and an
\emph{indirect} component arising from general-equilibrium adjustments of wages
and interest rates (and feeding into $\chi$). The direct component is 0.31 out
of the 0.68 excess return on fiscal spending ($\mu-1$) under heterogeneous
internalization, compared to 0.08 in the separable benchmark. This confirms the
central analytical prediction of Proposition~\ref{prop:xi-general}: \EC{}
introduces a new \emph{direct} transmission channel that operates through
household preferences rather than through income cyclicality.

\subsection{IRFs and transmission channels}
\label{subsec:irfs}

Figure~\ref{fig:irf} [to be generated: \path{code/hank/fig_irfs.py}] plots
impulse-response functions of output, private consumption, aggregate hours, and
the real wage to a 1\% persistent government-spending shock, under the four
configurations of Table~\ref{tab:multipliers}. The characteristic BoE-style
``shaded bands'' in the figure represent 68\% confidence intervals from a
Bayesian estimation of the complementarity parameters (details in
Section~\ref{sec:empirical}).

Three features of the IRFs stand out.
\textbf{(i)}~Private consumption \emph{rises} on impact under heterogeneous
internalization, whereas it falls slightly in the separable benchmark under
balanced-budget finance. The positive co-movement of private consumption with
$g$ -- a key empirical regularity \citep{Ramey2011} -- is recovered at
empirically plausible parameter values.
\textbf{(ii)}~The wage response is \emph{larger} under complementarity
because $\epsilon^{H}_{cg}<0$ raises the marginal utility of private consumption
for \HtM{} households, which shifts out their labor supply and puts upward
pressure on wages. This is the labor-supply channel of
Remark~\ref{rem:two-channels}.
\textbf{(iii)}~The multiplier decays at roughly the same rate across
configurations, reflecting the fact that $\delta$ (the THANK compounding
coefficient) depends on $\chi$ but not directly on $\theta^{j}$. The
differences in peak multipliers therefore propagate symmetrically through the
dynamic response.

\subsection{Sensitivity analysis}
\label{subsec:sensitivity}

Table~\ref{tab:sensitivity} [to be generated] reports impact multipliers for
the heterogeneous-internalization case across perturbations in each of the six
parameters the multiplier is most sensitive to: the \HtM{} share $\lambda$, the
saver persistence $s$, the Frisch elasticity $\phi$, the IES $\sigma$, the
estimated $\theta^{H}$, and the estimated $\theta^{S}$. The key findings:

\begin{itemize}[leftmargin=1.4em]
  \item The multiplier is most sensitive to $\theta^{H}$: a one-standard-error
    tighter bound on $\theta^{H}$ (i.e., moving from $-0.62$ to $-0.90$) raises
    $\mu$ from $1.60$ to $1.84$.
  \item The multiplier is convex in $\lambda$: amplification is strongest
    precisely when $\lambda$ is large -- consistent with
    Remark~\ref{rem:ineq-amp}.
  \item The multiplier is relatively insensitive to the IES $\sigma$ and
    Frisch $\phi$ along plausible ranges.
  \item Substitutability ($\theta^{S}>0$) dampens the multiplier only mildly
    under the baseline HtM share, confirming that \HtM{} complementarity is the
    dominant force in Corollary~\ref{cor:het-explicit}.
\end{itemize}

\subsection{Connecting HANK back to THANK}
\label{subsec:hank-vs-thank}

The close correspondence between the HANK multiplier and the THANK multiplier
under symmetric preferences (last rows of Table~\ref{tab:multipliers}) validates
the analytical sufficient-statistic approach. Under heterogeneous
internalization, THANK under-states the HANK multiplier by about 0.14 log
points. This difference is attributable to two features that THANK abstracts
from: (i) the continuous distribution of MPCs across the wealth distribution,
and (ii) the gradient of $\theta^{j}$ by wealth position (rather than a
two-type step function). A researcher who wishes to retain analytical
tractability while approximating HANK can therefore (a) use the THANK
expressions for qualitative conclusions and (b) apply a 10--15\% markup to
convert the THANK multiplier to a HANK counterpart at standard calibrations.
